% Options for packages loaded elsewhere
% Options for packages loaded elsewhere
\PassOptionsToPackage{unicode}{hyperref}
\PassOptionsToPackage{hyphens}{url}
\PassOptionsToPackage{dvipsnames,svgnames,x11names}{xcolor}
%
\documentclass[
  11pt,
]{article}
\usepackage{xcolor}
\usepackage[margin=1in]{geometry}
\usepackage{amsmath,amssymb}
\setcounter{secnumdepth}{-\maxdimen} % remove section numbering
\usepackage{iftex}
\ifPDFTeX
  \usepackage[T1]{fontenc}
  \usepackage[utf8]{inputenc}
  \usepackage{textcomp} % provide euro and other symbols
\else % if luatex or xetex
  \usepackage{unicode-math} % this also loads fontspec
  \defaultfontfeatures{Scale=MatchLowercase}
  \defaultfontfeatures[\rmfamily]{Ligatures=TeX,Scale=1}
\fi
\usepackage{lmodern}
\ifPDFTeX\else
  % xetex/luatex font selection
\fi
% Use upquote if available, for straight quotes in verbatim environments
\IfFileExists{upquote.sty}{\usepackage{upquote}}{}
\IfFileExists{microtype.sty}{% use microtype if available
  \usepackage[]{microtype}
  \UseMicrotypeSet[protrusion]{basicmath} % disable protrusion for tt fonts
}{}
\makeatletter
\@ifundefined{KOMAClassName}{% if non-KOMA class
  \IfFileExists{parskip.sty}{%
    \usepackage{parskip}
  }{% else
    \setlength{\parindent}{0pt}
    \setlength{\parskip}{6pt plus 2pt minus 1pt}}
}{% if KOMA class
  \KOMAoptions{parskip=half}}
\makeatother
% Make \paragraph and \subparagraph free-standing
\makeatletter
\ifx\paragraph\undefined\else
  \let\oldparagraph\paragraph
  \renewcommand{\paragraph}{
    \@ifstar
      \xxxParagraphStar
      \xxxParagraphNoStar
  }
  \newcommand{\xxxParagraphStar}[1]{\oldparagraph*{#1}\mbox{}}
  \newcommand{\xxxParagraphNoStar}[1]{\oldparagraph{#1}\mbox{}}
\fi
\ifx\subparagraph\undefined\else
  \let\oldsubparagraph\subparagraph
  \renewcommand{\subparagraph}{
    \@ifstar
      \xxxSubParagraphStar
      \xxxSubParagraphNoStar
  }
  \newcommand{\xxxSubParagraphStar}[1]{\oldsubparagraph*{#1}\mbox{}}
  \newcommand{\xxxSubParagraphNoStar}[1]{\oldsubparagraph{#1}\mbox{}}
\fi
\makeatother


\usepackage{longtable,booktabs,array}
\usepackage{calc} % for calculating minipage widths
% Correct order of tables after \paragraph or \subparagraph
\usepackage{etoolbox}
\makeatletter
\patchcmd\longtable{\par}{\if@noskipsec\mbox{}\fi\par}{}{}
\makeatother
% Allow footnotes in longtable head/foot
\IfFileExists{footnotehyper.sty}{\usepackage{footnotehyper}}{\usepackage{footnote}}
\makesavenoteenv{longtable}
\usepackage{graphicx}
\makeatletter
\newsavebox\pandoc@box
\newcommand*\pandocbounded[1]{% scales image to fit in text height/width
  \sbox\pandoc@box{#1}%
  \Gscale@div\@tempa{\textheight}{\dimexpr\ht\pandoc@box+\dp\pandoc@box\relax}%
  \Gscale@div\@tempb{\linewidth}{\wd\pandoc@box}%
  \ifdim\@tempb\p@<\@tempa\p@\let\@tempa\@tempb\fi% select the smaller of both
  \ifdim\@tempa\p@<\p@\scalebox{\@tempa}{\usebox\pandoc@box}%
  \else\usebox{\pandoc@box}%
  \fi%
}
% Set default figure placement to htbp
\def\fps@figure{htbp}
\makeatother





\setlength{\emergencystretch}{3em} % prevent overfull lines

\providecommand{\tightlist}{%
  \setlength{\itemsep}{0pt}\setlength{\parskip}{0pt}}



 


% ============================================================
%  Econ 502 — Problem Set / Handout Theme
% ============================================================

% ---------- packages ----------
\usepackage{enumitem}
\usepackage{tikz}
\usepackage{fancyhdr}
\usepackage{titlesec}
\usepackage{tcolorbox}
\usepackage{booktabs}

\tcbuselibrary{skins, breakable}

% ---------- colors ----------
\definecolor{csufblue}{HTML}{1E5288}
\definecolor{csufnavy}{HTML}{00274C}
\definecolor{accentteal}{HTML}{16A085}
\definecolor{lightgray}{HTML}{F5F5F5}
\definecolor{medgray}{HTML}{E0E0E0}
\definecolor{darktext}{HTML}{1A1A1A}

% ---------- page style ----------
\pagestyle{fancy}
\fancyhf{}
\renewcommand{\headrulewidth}{0.8pt}
\renewcommand{\headrule}{\hbox to\headwidth{\color{csufblue}\leaders\hrule height \headrulewidth\hfill}}
\fancyhead[L]{\small\color{csufblue}\textbf{Econ 502: Advanced Microeconomics}}
\fancyhead[R]{\small\color{csufblue}\thepage}
\fancyfoot{}

% first page: colored footer only
\fancypagestyle{plain}{%
  \fancyhf{}
  \renewcommand{\headrulewidth}{0pt}
  \fancyfoot[C]{\small\color{csufblue}\thepage}
}

% ---------- title block ----------
% Hook into Pandoc's \maketitle: add ruled lines around the title
\makeatletter
\renewcommand{\maketitle}{%
  \thispagestyle{plain}
  \begin{center}
    {\color{csufblue}\rule{\textwidth}{1.5pt}}\par
    \vspace{10pt}
    {\LARGE\bfseries\color{csufnavy}\@title\par}
    \vspace{8pt}
    {\color{csufblue}\rule{\textwidth}{1.5pt}}
  \end{center}
  \vspace{12pt}
}
\makeatother

% ---------- section headings ----------
\titleformat{\section}
  {\Large\bfseries\color{csufnavy}}
  {}{0pt}{}
  [\vspace{2pt}{\color{csufblue}\titlerule[0.8pt]}]

\titleformat{\subsection}
  {\large\bfseries\color{csufblue}}
  {}{0pt}{}

\titleformat{\subsubsection}
  {\normalsize\bfseries\color{accentteal}}
  {}{0pt}{}

\titlespacing*{\section}{0pt}{16pt}{8pt}
\titlespacing*{\subsection}{0pt}{12pt}{4pt}
\titlespacing*{\subsubsection}{0pt}{8pt}{4pt}

% ---------- tcolorbox environments ----------

% Answer/solution box (blue)
\newtcolorbox{answer}{
  colback   = csufblue!5,
  colframe  = csufblue,
  arc       = 3pt,
  boxrule   = 0.8pt,
  left      = 8pt,
  right     = 8pt,
  top       = 6pt,
  bottom    = 6pt,
  breakable,
}

% Key result box (teal, with optional title)
\newtcolorbox{keyresult}[1][]{
  colback      = accentteal!5,
  colframe     = accentteal,
  arc          = 3pt,
  boxrule      = 0.8pt,
  left         = 8pt,
  right        = 8pt,
  top          = 6pt,
  bottom       = 6pt,
  fonttitle    = \bfseries\color{white},
  coltitle     = white,
  colbacktitle = accentteal,
  title        = {#1},
  breakable,
}

% Note box (light gray)
\newtcolorbox{note}{
  colback  = lightgray,
  colframe = medgray,
  arc      = 3pt,
  boxrule  = 0.6pt,
  left     = 8pt,
  right    = 8pt,
  top      = 6pt,
  bottom   = 6pt,
  breakable,
}

% ---------- list styling ----------
\setlist[enumerate, 1]{label=(\alph*), leftmargin=*, topsep=4pt, itemsep=4pt}
\setlist[enumerate, 2]{label=(\roman*), leftmargin=*, topsep=2pt, itemsep=2pt}
\setlist[itemize]{leftmargin=*, topsep=4pt, itemsep=3pt}

% ---------- boxed answers ----------
% Colored \boxed command
\renewcommand{\boxed}[1]{\fcolorbox{csufblue}{csufblue!4}{$\displaystyle #1$}}

% ---------- hyperlink colors ----------
\AtBeginDocument{%
  \hypersetup{
    colorlinks = true,
    linkcolor  = csufblue,
    urlcolor   = accentteal,
    citecolor  = csufnavy,
  }%
}
\makeatletter
\@ifpackageloaded{caption}{}{\usepackage{caption}}
\AtBeginDocument{%
\ifdefined\contentsname
  \renewcommand*\contentsname{Table of contents}
\else
  \newcommand\contentsname{Table of contents}
\fi
\ifdefined\listfigurename
  \renewcommand*\listfigurename{List of Figures}
\else
  \newcommand\listfigurename{List of Figures}
\fi
\ifdefined\listtablename
  \renewcommand*\listtablename{List of Tables}
\else
  \newcommand\listtablename{List of Tables}
\fi
\ifdefined\figurename
  \renewcommand*\figurename{Figure}
\else
  \newcommand\figurename{Figure}
\fi
\ifdefined\tablename
  \renewcommand*\tablename{Table}
\else
  \newcommand\tablename{Table}
\fi
}
\@ifpackageloaded{float}{}{\usepackage{float}}
\floatstyle{ruled}
\@ifundefined{c@chapter}{\newfloat{codelisting}{h}{lop}}{\newfloat{codelisting}{h}{lop}[chapter]}
\floatname{codelisting}{Listing}
\newcommand*\listoflistings{\listof{codelisting}{List of Listings}}
\makeatother
\makeatletter
\makeatother
\makeatletter
\@ifpackageloaded{caption}{}{\usepackage{caption}}
\@ifpackageloaded{subcaption}{}{\usepackage{subcaption}}
\makeatother
\usepackage{bookmark}
\IfFileExists{xurl.sty}{\usepackage{xurl}}{} % add URL line breaks if available
\urlstyle{same}
\hypersetup{
  pdftitle={Problem Set 3},
  colorlinks=true,
  linkcolor={blue},
  filecolor={Maroon},
  citecolor={Blue},
  urlcolor={Blue},
  pdfcreator={LaTeX via pandoc}}


\title{Problem Set 3}
\usepackage{etoolbox}
\makeatletter
\providecommand{\subtitle}[1]{% add subtitle to \maketitle
  \apptocmd{\@title}{\par {\large #1 \par}}{}{}
}
\makeatother
\subtitle{Econ 502: Advanced Microeconomics}
\author{}
\date{}
\begin{document}
\maketitle


\subsection{Market Structure and
Welfare}\label{market-structure-and-welfare}

Consider a market with inverse demand \(P = 150 - Q\), where \(Q\) is
total output. There are \(n\) identical firms, each with constant
marginal cost \(c = 30\) and no fixed costs.

Denote each firm's output as \(q_i\) for \(i = 1, \ldots, n\), so that
total output is \(Q = \sum_{i=1}^n q_i\).

\subsubsection{Part I: Benchmarks}\label{part-i-benchmarks}

\begin{enumerate}
\def\labelenumi{\alph{enumi})}
\item
  \textbf{Perfect competition.} Suppose firms are price-takers
  (\(P = MC\)). Find the equilibrium price, total output, consumer
  surplus, and total surplus.
\item
  \textbf{Monopoly.} Suppose there is a single firm (\(n = 1\)). Derive
  the marginal revenue function. Find the profit-maximizing quantity,
  price, and profit. Calculate consumer surplus and the deadweight loss
  relative to perfect competition.
\end{enumerate}

\subsubsection{Part II: Price
Competition}\label{part-ii-price-competition}

\begin{enumerate}
\def\labelenumi{\alph{enumi})}
\setcounter{enumi}{2}
\tightlist
\item
  \textbf{Bertrand.} Suppose \(n \geq 2\) firms compete by
  simultaneously setting prices. Consumers buy from the cheapest firm;
  if prices are tied, they split the market equally. What is the Nash
  equilibrium? Compare the outcome to your answer in part (a) and
  briefly explain the economic logic.
\end{enumerate}

\subsubsection{Part III: Quantity
Competition}\label{part-iii-quantity-competition}

\begin{enumerate}
\def\labelenumi{\alph{enumi})}
\setcounter{enumi}{3}
\item
  \textbf{Cournot duopoly.} Now suppose \(n = 2\) firms compete by
  simultaneously choosing quantities, with total output
  \(Q = q_1 + q_2\). Write down firm 1's profit function and derive its
  best response function \(q_1^*(q_2)\). Find the Cournot-Nash
  equilibrium: each firm's output, the market price, and per-firm
  profit.
\item
  \textbf{Cournot with \(n\) firms.} Generalize part (d) to \(n\)
  identical firms. Show that in the symmetric Cournot-Nash equilibrium,
  per-firm output is given by:
\end{enumerate}

\[q^* = \frac{120}{n+1}\]

\begin{enumerate}
\def\labelenumi{\alph{enumi})}
\setcounter{enumi}{5}
\tightlist
\item
  \textbf{Convergence.} Using the result from part (e), fill in the
  following table and compute the markup \(\mu\), which is defined price
  relative to marginal cost: \[ \mu = \frac{P^*}{c} \]
\end{enumerate}

\begin{longtable}[]{@{}
  >{\centering\arraybackslash}p{(\linewidth - 8\tabcolsep) * \real{0.2000}}
  >{\centering\arraybackslash}p{(\linewidth - 8\tabcolsep) * \real{0.2000}}
  >{\centering\arraybackslash}p{(\linewidth - 8\tabcolsep) * \real{0.2000}}
  >{\centering\arraybackslash}p{(\linewidth - 8\tabcolsep) * \real{0.2000}}
  >{\centering\arraybackslash}p{(\linewidth - 8\tabcolsep) * \real{0.2000}}@{}}
\toprule\noalign{}
\begin{minipage}[b]{\linewidth}\centering
\(n\)
\end{minipage} & \begin{minipage}[b]{\linewidth}\centering
Per-firm output (\(q^*\))
\end{minipage} & \begin{minipage}[b]{\linewidth}\centering
Price (\(P^*\))
\end{minipage} & \begin{minipage}[b]{\linewidth}\centering
Per-firm profit (\(\pi^*\))
\end{minipage} & \begin{minipage}[b]{\linewidth}\centering
Markup (\(\mu\))
\end{minipage} \\
\midrule\noalign{}
\endhead
\bottomrule\noalign{}
\endlastfoot
1 & & & & \\
2 & & & & \\
5 & & & & \\
10 & & & & \\
50 & & & & \\
\end{longtable}

What happens to the markup as \(n\) gets large?

\subsubsection{Part IV: Welfare}\label{part-iv-welfare}

\begin{enumerate}
\def\labelenumi{\alph{enumi})}
\setcounter{enumi}{6}
\tightlist
\item
  \textbf{Deadweight loss.} Show that the deadweight loss from Cournot
  competition with \(n\) firms is given by:
\end{enumerate}

\[DWL = \frac{7{,}200}{(n+1)^2}\]

How does the DWL change as \(n\) increases? Compare the monopoly DWL
(\(n = 1\)) to the Cournot duopoly DWL (\(n = 2\)).

\begin{center}\rule{0.5\linewidth}{0.5pt}\end{center}




\end{document}
